%% thesis.tex 2014/04/11
%
% Based on sample files of unknown authorship.
%
% The Current Maintainer of this work is Paul Vojta.

\documentclass[masters]{ucbthesis}
\usepackage{amsmath}
\usepackage{latexsym}
\usepackage{amsfonts}
\usepackage[normalem]{ulem}
\usepackage{soul}
\usepackage{array}
\usepackage{amssymb}
\usepackage{extarrows}
\usepackage{graphicx}

\usepackage[utf8]{inputenc}
\usepackage[backend=biber,
natbib=true,
style=authoryear,
sorting=none,
isbn=true,
doi=true,
url=true,
]{biblatex}

\makeatletter
% special thanks to this guy
\newrobustcmd*{\parentexttrack}[1]{%
  \begingroup
  \blx@blxinit
  \blx@setsfcodes
  \blx@bibopenparen#1\blx@bibcloseparen
  \endgroup}

\AtEveryCite{%
  \let\parentext=\parentexttrack%
  \let\bibopenparen=\bibopenbracket%
  \let\bibcloseparen=\bibclosebracket}

\makeatother
\addbibresource{library.bib}

\usepackage{booktabs}
\usepackage{siunitx}
\usepackage{subfig}
\usepackage{wrapfig}
\usepackage{wasysym}
\usepackage{enumitem}
\usepackage{adjustbox}
\usepackage{ragged2e}
\usepackage[svgnames,table]{xcolor}
\usepackage{tikz}
% csv files and plots
\usepackage{pgfplots}
\pgfplotsset{compat=1.18}
\usepackage{csvsimple-l3}
\usepackage{longtable}
\usepackage{changepage}
\usepackage{setspace}
\usepackage{hhline}
\usepackage{multicol}
\usepackage{tabto}
\usepackage{float}
\usepackage{multirow}
\usepackage{makecell}
\usepackage{fancyhdr}
\usetikzlibrary{shapes.symbols,shapes.geometric,shadows,arrows.meta,
positioning,fit,backgrounds}
\tikzset{>={Latex[width=1.5mm,length=2mm]}}
\usepackage{tikz}
\usepackage[T1]{fontenc}
\usepackage{listings}

\urlstyle{same}

\renewcommand{\_}{\kern-1.5pt\textunderscore\kern-1.5pt}

%https://tex.stackexchange.com/questions/116534/lstlisting-line-wrapping
% I copied code from here to format things properly
% Credit to: Henri Menke who answered on StackExchange
\lstset{
  basicstyle=\ttfamily,
  columns=fullflexible,
  frame=single,
  breaklines=true,
  postbreak=\mbox{\textcolor{red}{$\hookrightarrow$}\space},
}

%%  Import package so that we can split the tex files into subfolders
% for easier compile and real-time edit

\usepackage{import}

%%%%%%%%%%%%%% for transfer function controllers %%%%%%%%%%%%%%
\usepackage{tikz}
\usetikzlibrary{shapes,arrows}
%%%% set citation style to square brackets


\begin{document}

% Declarations for Front Matter

\title{Derivations for Tampines Steam Tables and TAMPINES}
\author{Tom Lifesaver}
%\degreesemester{Spring}
%\degreeyear{1995}
%\degree{Doctor of Philosophy}
%\chair{Professor Richard Francis Sony}
%\othermembers{Professor Roger Spam \\
%  Associate Professor Michael Chex}
%% For a co-chair who is subordinate to the \chair listed above
%% \cochair{Professor Benedict Francis Pope}
%% For two co-chairs of equal standing (do not use \chair with this one)
%% \cochairs{Professor Richard Francis Sony}{Professor Benedict Francis Pope}
%\numberofmembers{3}
%% Previous degrees are no longer to be listed on the title page.
%% \prevdegrees{B.A. (University of Northern South Dakota at Hoople) 1978 \\
%%   M.S. (Ed's School of Quantum Mechanics and Muffler Repair) 1989}
%\field{Mathematics}
%% Designated Emphasis -- this is optional, and rare
%% \emphasis{Colloidal Telemetry}
%% This is optional, and rare
%% \jointinstitution{University of Western Maryland}
%% This is optional (default is Berkeley)
%% \campus{Berkeley}
%
%% For a masters thesis, replace the above \documentclass line with
%% \documentclass[masters]{ucbthesis}
%% This affects the title and approval pages, which by default calls this
%% document a "dissertation", not a "thesis".
%
%\maketitle
%% Delete (or comment out) the \approvalpage line for the final version.
%\approvalpage
%\copyrightpage

\include{abstract}

\begin{frontmatter}


% You can delete the \clearpage lines if you don't want these to start on
% separate pages.

\tableofcontents
\clearpage
\listoffigures
\clearpage
\listoftables

\begin{acknowledgements}
\end{acknowledgements}

\end{frontmatter}

\pagestyle{headings}

% (Optional) \part{First Part}

\part{Introduction}

\part{What is TAMPINES}

\part{Steam Tables}

\part{Rankine Cycle Derivations}

In this part, we consider a four control volume system representing a 
simple Rankine cycle. This will form the theoretical basis to 

\chapter{Mass, Energy and Momentum Balance}

In OpenFOAM, the PIMPLE algorithm is often used to solve for incompressible 
fluid flow \citep{devolder2015review}, this is one of the many standard 
solvers for solving flow in OpenFOAM \citep{penttinen2011pimplefoam}.
Jasak has published this in his PhD thesis, available online,
\citep{jasak1996error} on page 143-149.

\section{Raw forms of mass, momentum and energy balance}

Let's consider mass balance and momentum balance:
\begin{equation}
	\frac{\partial \dot{m}}{\partial t} = \sum_i m_i
\end{equation}


\begin{equation}
	\frac{\partial (\rho \textbf{u})}{\partial t} 
	+ \nabla \cdot (\rho \textbf{u} \textbf{u})
	= - \nabla P
	+ \nabla \cdot \tau 
	+ \rho g
\end{equation}

$\nabla \cdot \tau$ represents shear, $\rho g$ represents the gravity 
body force, $- \nabla P$ represents the pressure gradient.

\begin{equation}
	V_j \frac{\partial \rho_i h_i}{\partial t} + \sum_j \dot{m}_j h_j
	= Q_i - W_i 
	+ \sum_j \mathbf{\tau}_j \cdot \mathbf{u}_j
\end{equation}

Where $+ \sum_j \mathbf{\tau}_j \cdot \mathbf{u}_j $ is basically the 
friction loss terms where energy is transformed from kinetic energy 
to sensible heat. Q and W represent the thermodynamic work

\section{converting forms to mass flowrate}

It is often more convenient to convert the momentum equations to be 
in terms of mass flowrates. This is for 1D case.

To do so, we transform this by consider this 3D equation.
\begin{equation}
	\frac{\partial (\rho \textbf{u})}{\partial t} 
	+ \nabla \cdot (\rho \textbf{u} \textbf{u})
	= - \nabla P
	+ \nabla \cdot \tau 
	+ \rho g
\end{equation}

We need to consider just the x component, for a simple 1D case.

Then multiply by the cross sectional area.

\section{Approach to Solving nonlinear Equations}

Jasak in his PhD thesis outlines the methodology for solving this set of 
equation \citep{jasak1996error} on page 143-149. These were the foundations 
for what would eventually become OpenFOAM \cite{jasak2007openfoam} licensed 
under GPLv3.

The momentum equation is quadratic in u, which makes the equation solution 
nonlinear. Solving this nonlinear equation would either need a nonlinear 
solver such as Newton Raphson method, or solving a linearised form of 
these equations, and iterating to convergence. The latter was chosen 
\citep{jasak1996error}.

\begin{equation}
	\nabla \bullet (\mathbf{u u}) 
	= \sum_f \mathbf{S} \cdot \mathbf{u}_f\mathbf{u}_f
	= \sum_f F \mathbf{u}_f
	= a_p \mathbf{u}_p + \sum_N a_N \textbf{u}_n
\end{equation}

Section 3.3.1 within Jasak's thesis explains these terms.

First, we need to understand that Gauss's theorem is used in Finite 
Volume analysis, which is the bedrock of OpenFOAM solvers. 
$\sum_f \mathbf{S} \cdot \mathbf{u}_f\mathbf{u}_f$ represent the surface 
vector $\mathbf{S}$ and this will have a dot product with the velocity 
at the faces of the control volume $\mathbf{u}_f$.

Note that the subscript f represents values at the faces.

$\sum_f \mathbf{S} \cdot \mathbf{u}_f\mathbf{u}_f$
can be written in shorthand $\sum_f F \mathbf{u}_f$.

Now, we may not have the values of the velocities or fluxes at the 
faces. So we have to estimate it based on the properties at the 
centroid of the cell denoted $p$, and those of the neighbours denoted $n$.

Thus, the surface values denoted $f$ are estimated based on the values 
at the cell centroids $p$ and $n$ based on some kind of interpolation 
scheme. This can be linear interpolation, or some upwind differencing scheme.
This is described in page 81 of Jasak's thesis \citep{jasak1996error}.





\printbibliography

% \appendix
% \chapter{More Monticello Candidates}

\end{document}


